% META: {"id": "1SPE-SUITES-CO-004", "chapitre": "1SPE-SUITES", "type_objet": "corrige", "exercice_ref": "1SPE-SUITES-EX-004", "capacites_codes": ["C1"], "sources_inspiration": [], "status": "generated"}
% BEGIN-VERIFY
% from sympy import *
% n = symbols('n', integer=True, nonneg=True)
% u = n**2 - 2*n + 3
% assert u.subs(n, 0) == 3
% assert u.subs(n, 1) == 2
% assert u.subs(n, 2) == 3
% assert u.subs(n, 3) == 6
% assert u.subs(n, 5) == 18
% # Q3 : Amira a tort, u1 = 2 != u2 = 3
% assert u.subs(n, 1) != u.subs(n, 2)
% END-VERIFY
\begin{corrige}{1SPE-SUITES-EX-004}

\commentaireMarge{Calcul direct par substitution dans la formule explicite $u_n = n^2 - 2n + 3$.}

\textbf{Question 1.} On substitue les valeurs de $n$ dans $u_n = n^2 - 2n + 3$.

\[
u_0 = 0^2 - 2 \times 0 + 3 = 0 - 0 + 3 = 3.
\]
\[
u_1 = 1^2 - 2 \times 1 + 3 = 1 - 2 + 3 = 2.
\]
\[
u_2 = 2^2 - 2 \times 2 + 3 = 4 - 4 + 3 = 3.
\]
\[
u_3 = 3^2 - 2 \times 3 + 3 = 9 - 6 + 3 = 6.
\]

Donc $u_0 = 3$, $u_1 = 2$, $u_2 = 3$ et $u_3 = 6$.

\textbf{Question 2.} Pour $n = 5$ :
\[
u_5 = 5^2 - 2 \times 5 + 3 = 25 - 10 + 3 = 18.
\]

Donc $u_5 = 18$.

\commentaireMarge{On vérifie l'affirmation d'Amira par le calcul direct.}

\textbf{Question 3.} D'après les calculs de la question 1, on a $u_1 = 2$ et $u_2 = 3$.

Comme $2 \neq 3$, on a $u_1 \neq u_2$.

Donc Amira a \textbf{tort} : les termes $u_1$ et $u_2$ ne sont pas égaux.

\end{corrige}
